% GENERATED FILE. Edit canonical lessons, then run npm run generate:books.
% canonical-source-hash: fnv1a64:e35f8efa26648454
% canonical-lessons: PT-C466-chuviscar, PT-C466-desanuviar, PT-C466-poca, PT-C466-encharcado, PT-C466-gabardina, PT-R466-revisao-465-466, PT-R466-revisao-final-466

\chapter{Depois da chuva}
\label{ch:pt-depois-da-chuva}
\hlchaptermodality{466}

\begin{chapteropening}
\textbf{By the end of this chapter:} \emph{I can talk about rainy weather and what follows: drizzle, a sky that clears, puddles, getting soaked and the raincoat.}

Complete the last lesson of chapter 466: I can talk about rainy weather and what follows: drizzle, a sky that clears, puddles, getting soaked and the raincoat.
\end{chapteropening}

\section[chuviscar]{chuviscar — to drizzle}

\label{lesson:PT-C466-chuviscar}



\begin{quote}
Before the new one: say the Portuguese for the hurricane, then the Portuguese for the blizzard.
\end{quote}

\subsection*{You'll want to know: chuviscar}
\textbf{chuviscar} — \textquotedblleft{}to drizzle\textquotedblright{}.

\begin{quote}
Está a chuviscar, mas não preciso de guarda-chuva.
\end{quote}

\begin{quote}
\emph{It is drizzling, but I do not need an umbrella.}
\end{quote}

\emph{O chuvisco} — drizzle.

\begin{cousinweb}
\emph{chuviscar} is built on \emph{chuva} 'rain', which is inherited from Latin \emph{pluvia} (EtymDB).
\end{cousinweb}

\begin{grammarlens}[title={What you've built}]
One more word for telling, explaining and describing at B1.
\end{grammarlens}

\subsection*{Your turn}
\begin{itemize}
\raggedright
  \item \emph{Say it:} \emph{chuviscar}
  \item \emph{Say it:} \emph{chuviscar}, once more
  \item \emph{Say it:} a sentence of your own with \emph{chuviscar}
  \item \emph{Recall:} say the Portuguese for the hurricane, then the Portuguese for the blizzard, then say \emph{chuviscar} again
\end{itemize}

\begin{tcolorbox}[breakable,colback=teal!4,colframe=teal!35!black,title={Before you move on}]
What does chuviscar mean? (\textquotedblleft{}To drizzle\textquotedblright{}.) Say it once more.
\end{tcolorbox}
\section[desanuviar]{desanuviar — to clear up (of the sky); to relax}

\label{lesson:PT-C466-desanuviar}



\begin{quote}
Before the new one: say the Portuguese for the blizzard, then the Portuguese for to drizzle.
\end{quote}

\subsection*{You'll want to know: desanuviar}
\textbf{desanuviar} — \textquotedblleft{}to clear up (of the sky); to relax\textquotedblright{}.

\begin{quote}
À tarde, o céu desanuviou.
\end{quote}

\begin{quote}
\emph{In the afternoon, the sky cleared.}
\end{quote}

\emph{Desanuviar a cabeça} — to clear your head.

\begin{cousinweb}
\emph{desanuviar} is built from \emph{des-} and \emph{anuviar} 'to cloud over', on \emph{nuvem} 'cloud', which goes back to Latin \emph{nūbēs} (EtymDB).
\end{cousinweb}

\begin{grammarlens}[title={What you've built}]
One more word for telling, explaining and describing at B1.
\end{grammarlens}

\subsection*{Your turn}
\begin{itemize}
\raggedright
  \item \emph{Say it:} \emph{desanuviar}
  \item \emph{Say it:} \emph{desanuviar}, once more
  \item \emph{Say it:} a sentence of your own with \emph{desanuviar}
  \item \emph{Recall:} say the Portuguese for the blizzard, then the Portuguese for to drizzle, then say \emph{desanuviar} again
\end{itemize}

\begin{tcolorbox}[breakable,colback=teal!4,colframe=teal!35!black,title={Before you move on}]
What does desanuviar mean? (\textquotedblleft{}To clear up (of the sky); to relax\textquotedblright{}.) Say it once more.
\end{tcolorbox}
\section[a poça]{a poça — the puddle}

\label{lesson:PT-C466-poca}



\begin{quote}
Before the new one: say the Portuguese for to drizzle, then the Portuguese for to clear up.
\end{quote}

\subsection*{You'll want to know: a poça}
\textbf{a poça} — \textquotedblleft{}the puddle\textquotedblright{}.

\begin{quote}
O menino salta para as poças.
\end{quote}

\begin{quote}
\emph{The boy jumps into the puddles.}
\end{quote}

\emph{Uma poça de água} — a puddle of water.

\begin{cousinweb}
\emph{poça} is built on \emph{poço} 'well, pit' (EtymDB), which is inherited from Latin \emph{puteus}.
\end{cousinweb}

\begin{grammarlens}[title={What you've built}]
One more word for telling, explaining and describing at B1.
\end{grammarlens}

\subsection*{Your turn}
\begin{itemize}
\raggedright
  \item \emph{Say it:} \emph{a poça}
  \item \emph{Say it:} \emph{a poça}, once more
  \item \emph{Say it:} a sentence of your own with \emph{a poça}
  \item \emph{Recall:} say the Portuguese for to drizzle, then the Portuguese for to clear up, then say \emph{a poça} again
\end{itemize}

\begin{tcolorbox}[breakable,colback=teal!4,colframe=teal!35!black,title={Before you move on}]
What does a poça mean? (\textquotedblleft{}The puddle\textquotedblright{}.) Say it once more.
\end{tcolorbox}
\section[encharcado]{encharcado — soaked, drenched}

\label{lesson:PT-C466-encharcado}



\begin{quote}
Before the new one: say the Portuguese for to clear up, then the Portuguese for the puddle.
\end{quote}

\subsection*{You'll want to know: encharcado}
\textbf{encharcado} — \textquotedblleft{}soaked, drenched\textquotedblright{}.

\begin{quote}
Cheguei a casa encharcado.
\end{quote}

\begin{quote}
\emph{I got home soaked.}
\end{quote}

Built on \emph{o charco}, a pool of water.

\begin{cousinweb}
\emph{encharcado} is built from \emph{en-} and \emph{charco} 'pool, puddle', which was borrowed from Spanish \emph{charco} (EtymDB); the Spanish word's origin is not settled.
\end{cousinweb}

\begin{grammarlens}[title={What you've built}]
One more word for telling, explaining and describing at B1.
\end{grammarlens}

\subsection*{Your turn}
\begin{itemize}
\raggedright
  \item \emph{Say it:} \emph{encharcado}
  \item \emph{Say it:} \emph{encharcado}, once more
  \item \emph{Say it:} a sentence of your own with \emph{encharcado}
  \item \emph{Recall:} say the Portuguese for to clear up, then the Portuguese for the puddle, then say \emph{encharcado} again
\end{itemize}

\begin{tcolorbox}[breakable,colback=teal!4,colframe=teal!35!black,title={Before you move on}]
What does encharcado mean? (\textquotedblleft{}Soaked, drenched\textquotedblright{}.) Say it once more.
\end{tcolorbox}
\section[a gabardina]{a gabardina — the raincoat}

\label{lesson:PT-C466-gabardina}



\begin{quote}
Before the new one: say the Portuguese for the puddle, then the Portuguese for soaked.
\end{quote}

\subsection*{You'll want to know: a gabardina}
\textbf{a gabardina} — \textquotedblleft{}the raincoat\textquotedblright{}.

\begin{quote}
Saí de gabardina porque estava a chover.
\end{quote}

\begin{quote}
\emph{I went out in my raincoat because it was raining.}
\end{quote}

In Brazil: \emph{a capa de chuva}.

\begin{cousinweb}
\emph{gabardina} is usually traced to Spanish \emph{gabardina} 'long cloak'; its earlier origin is not settled, and EtymDB has no entry for it.
\end{cousinweb}

\begin{grammarlens}[title={What you've built}]
One more word for telling, explaining and describing at B1.
\end{grammarlens}

\subsection*{Your turn}
\begin{itemize}
\raggedright
  \item \emph{Say it:} \emph{a gabardina}
  \item \emph{Say it:} \emph{a gabardina}, once more
  \item \emph{Say it:} a sentence of your own with \emph{a gabardina}
  \item \emph{Recall:} say the Portuguese for the puddle, then the Portuguese for soaked, then say \emph{a gabardina} again
\end{itemize}

\begin{tcolorbox}[breakable,colback=teal!4,colframe=teal!35!black,title={Before you move on}]
What does a gabardina mean? (\textquotedblleft{}The raincoat\textquotedblright{}.) Say it once more.
\end{tcolorbox}
\section[(dialogue)]{Cidade, lojas, natureza e tempo}

\label{lesson:PT-R466-revisao-465-466}



\begin{quote}
Say the Portuguese for the gust and to blow.
\end{quote}

\subsection*{Your turn}
Say these aloud:

\begin{itemize}
\raggedright
  \item the gust, to blow, the storm, the hurricane, the blizzard
  \item to drizzle, to clear up, the puddle, soaked, the raincoat
\end{itemize}

\begin{tcolorbox}[breakable,colback=teal!4,colframe=teal!35!black,title={Before you move on}]
The gust? (\textbf{a rajada}.) To blow? (\textbf{soprar}.) The storm? (\textbf{o temporal}.) The hurricane? (\textbf{o furacão}.) The blizzard? (\textbf{a nevasca}.) To drizzle? (\textbf{chuviscar}.) To clear up? (\textbf{desanuviar}.) The puddle? (\textbf{a poça}.) Soaked? (\textbf{encharcado}.) The raincoat? (\textbf{a gabardina}.)
\end{tcolorbox}
\section[(dialogue)]{Depois da chuva — once more}

\label{lesson:PT-R466-revisao-final-466}



\begin{quote}
Say the Portuguese for to drizzle and to clear up.
\end{quote}

\subsection*{Your turn}
\begin{itemize}
\raggedright
  \item \emph{Say it:} to drizzle, to clear up, the puddle, soaked, the raincoat
\end{itemize}

\begin{tcolorbox}[breakable,colback=teal!4,colframe=teal!35!black,title={Before you move on}]
To drizzle? (\textbf{chuviscar}.) To clear up? (\textbf{desanuviar}.) The puddle? (\textbf{a poça}.) Soaked? (\textbf{encharcado}.) The raincoat? (\textbf{a gabardina}.)
\end{tcolorbox}
